\documentclass[12pt]{article}
\usepackage[margin=1in]{geometry}
\usepackage{setspace}
\usepackage{natbib}
\usepackage{hyperref}
\hypersetup{colorlinks=true, citecolor=blue, linkcolor=blue, urlcolor=blue}
\newcommand{\doi}[1]{\href{https://doi.org/#1}{\texttt{doi:#1}}}
\onehalfspacing
\title{Related Literature: A Capital Shock and Chilean Manufacturing\\
After the 2010 Maule Earthquake}
\author{Emily Liu}
\date{\today}

\begin{document}
\maketitle

\noindent This review situates the project within four literatures it draws on and one it
borrows its inference tools from. The organizing question is narrow: when a large, plausibly
exogenous shock destroys and then rebuilds a region's physical capital, which margin of the
firm actually moves, and who captures the returns to reconstruction? Each strand below answers
part of that question; none answers all of it at the establishment level, which is where this
project sits.

\section{Disasters and the trajectory of output}

The first and largest strand measures what disasters do to aggregate or regional economic
activity. The modern consensus is that large earthquakes leave persistent negative output
footprints rather than the ``creative destruction'' rebound that an accelerated--replacement
story might predict. \citet{aksoy2024}, assembling a global panel of earthquakes, find that
output falls and stays below trend for years, with the damage concentrated in lower-income
economies and those with weaker institutions; the capital-upgrading channel that could in
principle raise productivity is empirically dominated by the destruction of the existing stock.
For Chile specifically, \citet{diaz2024} compare the 2010 Maule earthquake to New Zealand's
Canterbury sequence and document that Chilean regional value added remained below its
counterfactual path well after the event, while \citet{benguria2022} reaches a broadly similar
conclusion for national GDP. The contemporaneous damage assessment in \citet{eclac2010} fixes
the scale of the shock---losses near 18\% of one year's GDP, with destruction concentrated in
housing and physical infrastructure---and motivates treating the earthquake as a capital shock
rather than a demand or institutional shock.

What this strand establishes is the \emph{sign and persistence} of the aggregate response. What
it cannot establish, because it works with regional or national aggregates, is the
\emph{composition} of that response inside the firm: whether the recovery is driven by
rebuilding capital, by rehiring, by changing the skill mix, or by none of these. Opening that
black box is the point of using establishment microdata, and it is where this project departs
from the aggregate-trajectory literature it otherwise rests on.

\section{Disasters and local labor markets}

A smaller strand looks past aggregates to the local labor market. \citet{belasen2009} use
quarterly establishment data to show that Florida hurricanes raise earnings in directly hit
counties---consistent with a local labor-demand pull from reconstruction net of any supply
loss---while depressing them in neighboring counties. \citet{notowidigdo2020} provides the
organizing frame for thinking about who bears a \emph{local} demand shock at all: because less-skilled
workers are less geographically mobile but better insured by transfers, the welfare incidence of
a local shock is distributed differently across skill groups than a naive wage regression would
suggest. This is the literature closest in spirit to the project's wage analysis, and it supplies
two cautions that the empirical design has to take seriously. First, a wage response in the hit
region can reflect either a demand pull or a compositional change in who remains---so a labor-supply
(out-migration) confound must be ruled out before any wage movement is read as demand. Second,
the relevant skill groups may be integrated across a wider labor market than the treated region,
which blunts the local skill-premium response one might otherwise expect.

The project takes both cautions on board: it uses a household survey (CASEN) to test directly
whether the unskilled population fell in harder-hit regions (it did not), and it interprets the
fragility of the skill-premium estimate through the lens of geographic integration rather than
as evidence against complementarity. What this strand does not do---again for want of
capital data---is connect the labor-market response to what is happening on the capital side of
the same establishments. That connection is the project's contribution.

\section{Capital--skill complementarity and the skill premium}

The motivating hypothesis comes from the capital--skill complementarity literature.
\citet{krusell2000} show that if equipment capital is more complementary with skilled than with
unskilled labor, then the secular fall in equipment prices can account for much of the rise in
the skill premium. Read as a prediction about shocks, this says that a surge in equipment
capital should raise the relative demand for, and hence the relative wage of, skilled workers.
The earthquake offers a natural test: reconstruction is, mechanically, a pulse of new capital.

\citet{lewis2011} complicates the prediction in a way that turns out to matter for the results.
Using U.S.\ local labor markets, he shows that factor supplies shape technology adoption in
equilibrium---an inflow of less-skilled labor slows the adoption of automation machinery---so the
observed relationship between capital and the skill premium is an equilibrium object, not a
structural parameter that a reduced-form regression recovers. The implication for this project is
that a reduced-form ``Post~$\times$~intensity'' coefficient on the skill premium identifies a
\emph{net} incidence, not the structural complementarity parameter, and that a small or fragile
estimate is consistent with---rather than a refutation of---complementarity once local factor
supplies and cross-regional integration are allowed to respond. The project's finding (a
skill-premium response that does not survive the parallel-trends sensitivity analysis) is
reported in exactly these restrained terms.

\section{Incidence: who captures the returns to capital}

If reconstruction raises returns to capital, the natural question is who ends up with them. The
incidence literature has answered the analogous question for \emph{fiscal} shocks to the price of
capital and labor. \citet{suarezzidar2016} estimate that the burden of U.S.\ state corporate
taxes is split among firm owners, workers, and landowners, with a substantial share falling on
owners; \citet{fuest2018} find that roughly half of the German local business-tax burden is
borne by workers through lower wages; and \citet{gruber1997}, studying Chile's own payroll-tax
reductions, finds the change almost fully shifted into wages with no employment effect. These
papers share a design---a tax wedge that moves the cost of a factor---and a question---what share
of the change shows up in wages versus profits.

The 2010 earthquake is the same incidence question with a different instrument: the shock is
\emph{physical} rather than fiscal. It destroys capital and finances its replacement, and one can
ask what share of the reconstruction shows up as higher labor compensation versus as
rebuilt capital and recovering returns. The project answers this directly by decomposing the
post-shock response across investment, the wage bill, labor's share of output, and a profit
proxy on a common footing. The finding---investment and the total wage bill rise while labor's
share does not---places the incidence on capital, and is, to my knowledge, the first
establishment-level incidence decomposition of a physical capital-destruction shock. The fiscal
incidence literature supplies both the question and the accounting; the physical shock supplies
the variation.

\section{Inference: continuous-treatment difference-in-differences}

The project's identifying variation is continuous (ground-shaking intensity by region) and its
panel has few clusters (about fifteen regions) and a known pre-trend (central regions were
already growing faster). Three methodological papers govern how the results are reported.
\citet{roth2022} shows that pre-testing for parallel trends and conditioning on passing the test
distorts inference and has low power against exactly the trends that matter; the project therefore
does not treat a passed pre-trend test as a license. \citet{rambachanroth2023} provide the
constructive alternative used throughout: instead of assuming parallel trends, bound the
post-period bias by the magnitude of pre-period violations (relative-magnitudes and smoothness
restrictions) and report the breakdown value at which the result disappears. This is what
downgrades the investment effect from ``significant'' to ``robust but suggestive'' and what
collapses the skill-premium effect. Finally, \citet{cameron2008} supply the wild cluster
bootstrap used to keep inference honest with so few clusters. Together these tools are the reason
the paper's claims are stated conservatively rather than at face-value significance.

\section{Where this project sits}

Each literature answers part of the question and stops at the edge of its data. The
aggregate-trajectory papers establish that the shock is real and persistent but cannot see inside
the firm. The local-labor-market papers see wages and employment but not capital. The
complementarity papers predict a skill-premium response but, after \citet{lewis2011}, warn that a
reduced form will not recover it cleanly. The incidence papers supply the right accounting but
for fiscal, not physical, shocks. This project's contribution is to bring establishment-level
capital and labor data to bear on all four at once: it shows that the only margin that moves is
reconstruction investment, that output, employment, productivity, the skill premium, and labor's
share do not, and that the incidence of the rebuilding therefore falls on capital rather than
labor---reconstruction without reallocation. The methodological literature is what licenses
stating that conclusion with the appropriate caution.

\bibliographystyle{plainnat}
\begin{thebibliography}{99}
\bibitem[Aksoy et al.(2024)]{aksoy2024} Aksoy, C.~G., Chupilkin, M., Koczan, Z., \& Plekhanov, A.
(2024). Unearthing the economic and social consequences of earthquakes. \textit{IZA Discussion
Paper No.\ 17108}.
\bibitem[Belasen \& Polachek(2009)]{belasen2009} Belasen, A.~R., \& Polachek, S.~W. (2009). How
disasters affect local labor markets: The effects of hurricanes in Florida. \textit{Journal of
Human Resources}, 44(1), 251--276. \doi{10.3368/jhr.44.1.251}
\bibitem[Benguria(2022)]{benguria2022} Benguria Garibaldi, S. (2022). Shaky growth: Chile's
earthquake and its effect on GDP. Master's thesis, Stockholm School of Economics.
\bibitem[Cameron et al.(2008)]{cameron2008} Cameron, A.~C., Gelbach, J.~B., \& Miller, D.~L.
(2008). Bootstrap-based improvements for inference with clustered errors. \textit{Review of
Economics and Statistics}, 90(3), 414--427. \doi{10.1162/rest.90.3.414}
\bibitem[D\'iaz et al.(2024)]{diaz2024} D\'iaz, D., Paniagua, P., \& Larroulet, C. (2024).
Earthquakes and the wealth of nations: The cases of Chile and New Zealand. \textit{Working
paper}, arXiv:2405.12041.
\bibitem[ECLAC(2010)]{eclac2010} Economic Commission for Latin America and the Caribbean. (2010).
\textit{The Chilean Earthquake of 27 February 2010: An Overview}. Santiago: ECLAC.
\bibitem[Fuest et al.(2018)]{fuest2018} Fuest, C., Peichl, A., \& Siegloch, S. (2018). Do
higher corporate taxes reduce wages? \textit{American Economic Review}, 108(2), 393--418.
\doi{10.1257/aer.20130570}
\bibitem[Gruber(1997)]{gruber1997} Gruber, J. (1997). The incidence of payroll taxation:
Evidence from Chile. \textit{Journal of Labor Economics}, 15(S3), S72--S101. \doi{10.1086/209877}
\bibitem[Krusell et al.(2000)]{krusell2000} Krusell, P., Ohanian, L.~E., R\'ios-Rull, J.-V.,
\& Violante, G.~L. (2000). Capital-skill complementarity and inequality.
\textit{Econometrica}, 68(5), 1029--1053. \doi{10.1111/1468-0262.00150}
\bibitem[Lewis(2011)]{lewis2011} Lewis, E. (2011). Immigration, skill mix, and capital-skill
complementarity. \textit{Quarterly Journal of Economics}, 126(2), 1029--1069.
\doi{10.1093/qje/qjr011}
\bibitem[Notowidigdo(2020)]{notowidigdo2020} Notowidigdo, M.~J. (2020). The incidence of local
labor demand shocks. \textit{Journal of Labor Economics}, 38(3), 687--725. \doi{10.1086/706048}
\bibitem[Rambachan \& Roth(2023)]{rambachanroth2023} Rambachan, A., \& Roth, J. (2023). A more
credible approach to parallel trends. \textit{Review of Economic Studies}, 90(5), 2555--2591.
\doi{10.1093/restud/rdad018}
\bibitem[Roth(2022)]{roth2022} Roth, J. (2022). Pretest with caution: Event-study estimates
after testing for parallel trends. \textit{American Economic Review: Insights}, 4(3), 305--322.
\doi{10.1257/aeri.20210236}
\bibitem[Su\'arez Serrato \& Zidar(2016)]{suarezzidar2016} Su\'arez Serrato, J.~C., \& Zidar,
O. (2016). Who benefits from state corporate tax cuts? \textit{American Economic Review},
106(9), 2582--2624. \doi{10.1257/aer.20141702}
\end{thebibliography}

\end{document}
